\documentclass{tudelft}
\geometry{margin=0.5mm}
\usepackage{circuitikz}


\newcommand{\thecourse}{Lineare circuits}

% ---------- DOCUMENT ----------
\begin{document}
\twocolumn

% \title{
%     \includesvg[width=0.75\linewidth]{TU_Delft_Logo} \\[2em]
%     \thecourse: Sample circuits
% }
% \author{Simon Zweers}
% \date{\today}
% \degree{Electrical Engineering}
% \faculty{Electrical Engineering, Mathematics \& Computer Science}

% \maketitle
% \thispagestyle{fancy}
% \vspace{1em}
% \hrule
% \vspace{1em}

\newpage

% TOC bs
\setcounter{secnumdepth}{0}
% \tableofcontents
% \newpage

% Main Text
Cheatsheet \\
Current/charge: $\displaystyle i(t) \triangleq \frac{dq}{dt} \iff q=\int_{t_0}^{t}idt + q(t_0)$ \\ 
Voltage: $\displaystyle v_{ab} \triangleq \frac{dw}{dq} $ \\
Power: $\displaystyle p = \frac{dw}{dt} = \frac{dw}{dq}\cdot\frac{dq}{dt} = v\cdot i $ \\
$\displaystyle p = i^2R = \frac{v^2}{R}$ \\ 
Energy: $\displaystyle w(t) = \int_{t_0}^{t}p(t)dt = \int_{t_0}^{t}v(t)i(t)dt $\\ 
Resitance: $\displaystyle R=\rho\frac{\ell}{A} $ \\ 
Parallel resistance: $\displaystyle R_{eq} = \frac{R_1 R_2}{R_1 + R_2} $ \\
{\bf KCL}: $\displaystyle \sum_{j=1}^{n} i_j(t) = 0 $ \hspace{0.5cm}
{\bf KVL}: $\displaystyle \sum_{j=1}^{n} v_j(t) = 0 $ \\
% \include{circuit1.tex}
% \include{circuit2.tex}
% \include{circuit3.tex}
% \include{code1.tex}

% VOLTAGE DIVISION
\begin{center}
	\bf Voltage division:
\end{center}
\begin{minipage}{0.5\linewidth}
\begin{circuitikz}[american]
	\draw (0,4) to [V, l=$v_{in}$] (0,0)
	-- (3,0)
	to [R, l=$R_2$] (3,2)
	to [R, l=$R_1$] (3,4)
	-- (0,4)
	;
	\draw (3.5, 4) to [open, v=$V_1$] (3.5,2) ;
	\draw (3.5, 2) to [open, v=$V_2$] (3.5,0) ;
\end{circuitikz}
\end{minipage}\hfill
\begin{minipage}{0.45\linewidth}
$\displaystyle v_1 = v_{in}\frac{R_1}{R_1 + R_2}$

\vspace{0.5cm}
$\displaystyle v_2 = v_{in}\frac{R_2}{R_1 + R_2}$
\end{minipage}

%CURRENT DIVISION
\begin{center}
	\bf Current Division:
\end{center}
\begin{minipage}{0.5\linewidth}
	\begin{circuitikz}[american]
		\draw (0,0) to [I, l_=$i_{in}$] (0,3)
		-- (2,3) to [R, l=$R_1$, i=$i_1$] (2,0) -- (0,0)
		;
		\draw (2,3) -- (4,3)
		to [R, l_=$R_2$, i_=$i_2$] (4,0)
		-- (2,0)
		;
	\end{circuitikz}

\end{minipage}\hfill
\begin{minipage}{0.45\linewidth}
$\displaystyle i_1 = i_{in}\frac{R_2}{R_1 + R_2}$

\vspace{0.5cm}
$\displaystyle i_2 = i_{in}\frac{R_1}{R_1 + R_2}$
\end{minipage}

\end{document}
